\plannedchapter{操作系统噪声、时间源和性能实验边界}
{本章解释为什么测量不是简单计时：OS 调度、CPU 频率、虚拟化、page fault、interrupt 和后台任务都会改变结果。}
{\item wall time、CPU time、cycle counter 有什么差别？
\item 为什么第一次运行常常更慢？
\item WSL2、虚拟机和裸机 Linux 的测量边界在哪里？}
{\item scheduler、context switch、timer interrupt 的基本模型。
\item TSC、clock source、\code{std::chrono} 的抽象边界。
\item 频率 scaling、turbo、温度和功耗限制。}
{写一个重复测量程序，记录 min/median/max；比较首次访问和 warmup 后访问；记录环境信息并解释噪声来源。}
